% Declarations, AI statement and Open Practices Statement (BRM). The Open Practices Statement
% must sit immediately before the References, so it is last here. Headings follow the BRM
% submission guidelines (Declarations list; Open Practices section), checked 4 Oct 2026.

\section*{Declarations}

\paragraph{Funding.} No funding was received for this work.

\paragraph{Competing interests.} The author has no competing interests to declare.

\paragraph{Ethics approval.} Not applicable. The study analyzes public, de-identified NHANES microdata, released data from other studies and text generated by language models. No human participants were recruited.

\paragraph{Consent to participate.} Not applicable.

\paragraph{Consent for publication.} Not applicable.

\paragraph{Availability of data and materials.} The corpus, the persona registries, the original generation code for both framings, the outputs of every analysis and every receipt file are at \url{https://github.com/pskeough/A-Validity-Ladder-for-LLM-Simulated-Populations}. NHANES data are public from the National Center for Health Statistics, and the repository includes the code that extracts them. Raw files of the external releases are not redistributed, and the repository names the public source of each. SubPOP's training data, used only to check question overlap, derive from Pew Research Center's American Trends Panel and the General Social Survey \citep{davern2024gss}. The opinions expressed herein, including any implications for policy, are those of the author and not of the survey research centers.

\paragraph{Code availability.} The analysis scripts are in the same repository. Every level-2 external release can be read through \texttt{paper\_brm/external/harness}, whose self-test reproduces the verdicts of the article's own scripts on 14,349 stored contrasts with no mismatch, and the personality release was read through the article's level-1 and level-4 code. The Python package \texttt{validity-ladder} (version 0.1.1) implements the gate, levels 1 to 4 and step 0, taking the invariance verdict R4 as an input. It installs with \texttt{pip} and carries tests that reproduce the article's verdicts from the stored outputs.

\paragraph{Authors' contributions.} Patrick Keough is the sole author and takes responsibility for all content.

\paragraph{Generative AI disclosure.} Generative AI tools (Anthropic's Claude, through Claude Code) were used for language editing, for LaTeX formatting, and to write the Python analysis and verification scripts to the author's specification. The research questions, study design, statistical specification and interpretations are the author's own, and all generated code was reviewed by the author. The verification scripts that check every numeric claim in this manuscript against its receipt are released with the data.

\section*{Open Practices Statement}

Data, materials and analysis code are available at \url{https://github.com/pskeough/A-Validity-Ladder-for-LLM-Simulated-Populations}, release commit \texttt{COMMIT}, where \texttt{paper\_brm/manuscript/supplement/S2\_receipts.csv} maps every number in the article to the file that prints it. None of the analyses were preregistered.
